\Subhead{1.6.2}{Взаимная синхронизация колебаний двух генераторов}
Два самостоятельных генератора через связь воздействуют на фазы друг друга. При слабой связи, быстро восстанавливающихся амплитудах и синусоидальной фазовой характеристике без запаздывания:
\begin{equation}\label{eq:mutuallocking}
 \dot\theta_1=\omega_1+K_1\sin(\theta_2-\theta_1),\qquad
 \dot\theta_2=\omega_2+K_2\sin(\theta_1-\theta_2).
\end{equation}
$K_1$ --- скорость фазового воздействия второго генератора на первый, $K_2$ --- первого на второй; размерность с$^{-1}$. Это не безразмерный коэффициент взаимной индукции. Для $K_1,K_2>0$, $\psi=\theta_1-\theta_2$:
\begin{equation}
 \dot\psi=\omega_1-\omega_2-(K_1+K_2)\sin\psi.
\end{equation}
Устойчивая синхронизация существует при $|\omega_1-\omega_2|<K_1+K_2$, $\cos\psi_*>0$. Подставляя $\sin\psi_*=(\omega_1-\omega_2)/(K_1+K_2)$ в любое исходное уравнение, получаем общую частоту
\begin{equation}
 \Omega_{\rm sync}=\frac{K_2\omega_1+K_1\omega_2}{K_1+K_2}.
\end{equation}
Она ближе к частоте генератора, который меньше подвержен чужому воздействию. При $K_2\to0$ второй становится внешним эталоном, при симметричной связи частота равна полусумме. Вне полосы фаза проскальзывает, как в~\Qref{1.3.1}.

Резистивная связь обменивает ток, зависящий от разности напряжений; индуктивная или ёмкостная --- реактивную энергию. Реальная схема может давать фазовое запаздывание, косинусные члены и связь с амплитудой. Для чисто реактивной связи нельзя автоматически брать положительные $K_j$ из модели выше: нередко нужен амплитудно-фазовый расчёт. Сильная связь может подавлять генерацию; фазовая модель с уже ненулевыми амплитудами такой режим не описывает.
% Формулы фазовой модели выведены здесь. Привязка K_j к номиналам
% конкретной схемы курса не установлена; универсальная формула не заявляется.
